\documentclass[11pt,a4paper]{article}
\usepackage[T1]{fontenc}
\usepackage[utf8]{inputenc}
\usepackage{amsmath,amssymb,amsthm}
\usepackage[margin=1in]{geometry}
\usepackage[colorlinks=true,linkcolor=blue,urlcolor=blue]{hyperref}
\theoremstyle{plain}
\newtheorem{theorem}{Theorem}
\newtheorem{lemma}[theorem]{Lemma}
\newtheorem{proposition}[theorem]{Proposition}
\newtheorem{corollary}[theorem]{Corollary}
\theoremstyle{definition}
\newtheorem{remark}[theorem]{Remark}

\title{The empirical recurrence for OEIS A267669, proved by transfer matrix}
\author{Adrian Perez Fontelles\\ \small Independent researcher}
\date{3 September 2026}
\begin{document}
\maketitle

\begin{abstract}
OEIS A267669 counts arrays in which the REPEATED VALUES of each row, and of each column, stand in a
prescribed relation to one another in the order they occur. The entry carries an empirical
recurrence of order $21$ contributed by R.~H.~Hardin. It is true, and it is decidable
rather than empirical: one of the two scans runs inside a slice of the array and the other
across slices, and the second needs only the previous slice together with the last repeated
value so far in each column of the scan, so the count is a walk count on $S=1122$ states.
\end{abstract}

\noindent\small 2020 Mathematics Subject Classification. 05A15, 05A18, 15A18.\normalsize

\section{The conjecture}

OEIS A267669 is ``Number of 4Xn 0..1 arrays with every repeated value in every row unequal to the previous repeated value, and in every column equal to the previous repeated value, and new values introduced in row-major sequential order.''. It has offset $1$ and begins
\[
7, 98, 453, 3143, 10933, 60401, 188301, 951113,\ \dots
\]
The entry states, as a conjecture contributed by its author and never marked settled:
\begin{quote}
Empirical: a(n) = 5*a(n-1) +46*a(n-2) -255*a(n-3) -804*a(n-4) +5287*a(n-5) +6393*a(n-6) -57990*a(n-7) -17379*a(n-8) +368289*a(n-9) -70362*a(n-10) -1394997*a(n-11) +645844*a(n-12) +3138583*a(n-13) -1858369*a(n-14) -4077534*a(n-15) +2440404*a(n-16) +2940332*a(n-17) -1387260*a(n-18) -1138872*a(n-19) +275184*a(n-20) +183456*a(n-21)
\end{quote}
As of the ``Last modified'' line on the live entry (Jan 19 2016 08:10:24, revision 4) this is still
recorded as empirical, and nothing on the entry records it as proved.

\section{What is being counted}

Read a row of the array from left to right. A REPEATED VALUE is an entry equal to the entry
immediately before it; the row thus produces the sequence of values at the places where two
equal entries stand side by side. The entry asks that each repeated value of a row be
different from the previous repeated value of that row, and that each repeated value of a column,
read downwards, be equal to the previous repeated value of that column. The FIRST repeated
value of a scan has no predecessor and is unconstrained; a scan with no two equal neighbours
imposes nothing at all. The entry writes the array with $n$ as the number of COLUMNS, so the array grows one COLUMN at a time: the column scan is the one that runs inside a slice and the row scan is the one that runs across slices.

The entry closes with ``new values introduced in row-major sequential order'', so what is
counted is arrays up to renaming the letters. Both relations here are built from equality
alone, so the condition survives renaming and the usual inversion applies: with $N_j$ the
classes on exactly $j$ letters and $L_i$ the admissible arrays over an alphabet of $i$ letters,
\[
L_i=\sum_j N_j\,i(i-1)\cdots(i-j+1),\qquad
a=N_1+\dots+N_{2}=\bigl(\mathbf{1}^{\!\top}F^{-1}\bigr)L .
\]
An entry of this family whose relation compares letters by size, or does arithmetic on them,
carries no such clause and is counted over its own alphabet instead; the two cases are not
mixed.

\section{One scan inside a slice, one across}

Take the array a $column$ at a time, the direction in which it grows.

\begin{lemma}\label{lem:state}
The $column$ scans are settled as each $column$ is written. The $row$ scans are
settled by the previous $column$ together with, for each position, the last repeated value
so far in that $row$.
\end{lemma}

\begin{proof}
A $column$ scan compares neighbours inside one $column$, so writing the
$column$ left to right and carrying the last repeated value seen decides it. In a
$row$ scan a repeat happens at the new $column$ exactly when the new entry equals the
entry in the same position of the previous $column$, and the only thing the relation needs
besides that entry is the previous repeated value of that $row$.
\end{proof}

So the state is the pair (previous $column$, vector of last repeated values), the second
component taking $1+2$ values at each position --- the letters together with a symbol for
``none yet''. Assembling the pieces for $i=1,\dots,2$ as a disjoint union, with $\iota$ carrying the weights above cleared of denominators,
\[
a(n)\;=\;\iota^{\!\top}M^{\,j}\tau ,\qquad S=1122,\quad\tau=\mathbf{1} .
\]
The walk index $j$ and the entry's own $n$ differ by a constant, read off by matching the
published terms rather than assumed. In particular $a$ is $C$-finite.

\section{The proof}

\begin{lemma}\label{lem:crit}
Let $q(t)=t^{21}-\sum_i c_it^{21-i}$ be the characteristic polynomial of the
conjectured recurrence. The residual $a(n)-\sum_ic_ia(n-i)$ equals
$\iota^{\!\top}M^{\,j}q(M)\tau$ for the corresponding walk index $j$, so the recurrence
holds from some point on if and only if $\iota^{\!\top}M^{j}q(M)\tau=0$ for all large $j$,
and it is enough to examine $j<S$.
\end{lemma}

\begin{proof}
Factor $M^{j}$ out of $M^{j+21}-\sum_ic_iM^{j+21-i}$. For the bound, $M^{S}$
is an integer combination of $I,M,\dots,M^{S-1}$ by Cayley--Hamilton, so the numbers
$u_j=\iota^{\!\top}M^{j}q(M)\tau$ obey the monic recurrence given by the characteristic
polynomial of $M$; $S$ consecutive zeros force every later one, and the last nonzero $u_j$
pins the threshold exactly.
\end{proof}

\begin{theorem}
\[
a(n)\;=\;5\,a(n-1) + 46\,a(n-2) - 255\,a(n-3) - 804\,a(n-4) + 5287\,a(n-5) + 6393\,a(n-6) - 57990\,a(n-7) - 17379\,a(n-8) + 368289\,a(n-9) - 70362\,a(n-10) - 1394997\,a(n-11) + 645844\,a(n-12) + 3138583\,a(n-13) - 1858369\,a(n-14) - 4077534\,a(n-15) + 2440404\,a(n-16) + 2940332\,a(n-17) - 1387260\,a(n-18) - 1138872\,a(n-19) + 275184\,a(n-20) + 183456\,a(n-21)
\]
for every $n>21$.
\end{theorem}

\begin{proof}
The vector $w=q(M)\tau$ was formed by $21$ matrix--vector products in exact integer
arithmetic and $\iota^{\!\top}M^{j}w$ evaluated for $j=0,1,\dots$, again exactly; those
integers vanish from the index corresponding to $n=21+1$ onwards and the run continued
until the working vector was identically zero.
\end{proof}

\section{Verification}

Four checks, each independent of the algebra above.

First, the model reproduces all $22$ terms the entry publishes, exactly, in integer
arithmetic.

Second, the count was recomputed for the smallest arrays straight from the definition:
enumerate every array of that size, form each row's and each column's list of repeated values
by scanning it, and test consecutive members of those lists against the relation. No state, no
window and no transfer matrix are involved, and the published terms came back.

Third, the conjectured recurrence was evaluated directly on the published terms, with no
matrices involved, and holds wherever the proved range and the data overlap.

Fourth, the annihilation test of Lemma~\ref{lem:crit} was carried out in exact integer
arithmetic on the whole vector rather than on a sampled prefix.

\begin{thebibliography}{9}
\bibitem{oeis} The OEIS Foundation, \emph{The On-Line Encyclopedia of Integer
Sequences}, \url{https://oeis.org}, 2026. Sequence A267669.
\bibitem{stanley} R.~P.~Stanley, \emph{Enumerative Combinatorics}, Volume 1, 2nd ed.,
Cambridge University Press, 2011. (Transfer-matrix method, Section 4.7.)
\bibitem{horn} R.~A.~Horn and C.~R.~Johnson, \emph{Matrix Analysis}, 2nd ed., Cambridge
University Press, 2013.
\end{thebibliography}

\end{document}
